\documentclass[
    language=english,
    doctype=study-guide,
    institution=none,
]{../../omnilatex}

\RequirePackage{config/document-settings}

\begin{document}

\maketitle

\section{Exam Overview}

\noindent\textbf{Course:} Linear Algebra \\
\textbf{Exam Date:} March 15, 2026 \\
\textbf{Format:} 50 multiple-choice (2 pts each) + 5 short-answer (10 pts each) \\
\textbf{Total:} 150 points \\
\textbf{Allowed Materials:} One handwritten $3 \times 5$ index card

\bigskip
\noindent\textbf{Topics Covered:}
\begin{enumerate}
    \item Systems of linear equations and Gaussian elimination
    \item Vector spaces and subspaces
    \item Linear transformations and matrix representations
    \item Eigenvalues, eigenvectors, and diagonalization
    \item Inner product spaces and orthogonality
\end{enumerate}

\section{Topic 1: Systems of Linear Equations}

\subsection{Key Concepts}

A system of linear equations can be written in matrix form as $A\mathbf{x} = \mathbf{b}$, where $A$ is the $m \times n$ coefficient matrix, $\mathbf{x}$ is the vector of unknowns, and $\mathbf{b}$ is the right-hand side.

\begin{itemize}
    \item \textbf{Row echelon form (REF):} All nonzero rows are above any all-zero rows; the leading coefficient of each nonzero row is to the right of the leading coefficient of the row above.
    \item \textbf{Reduced row echelon form (RREF):} REF plus each leading coefficient is 1 and is the only nonzero entry in its column.
    \item \textbf{Pivot variables:} Columns containing a leading 1 in RREF.
    \item \textbf{Free variables:} Columns not containing a leading 1.
\end{itemize}

\subsection{Key Formulas}

\[
\text{Number of solutions} =
\begin{cases}
    0 & \text{if a row of the augmented matrix is } [0 \;\cdots\; 0 \mid b] \text{ with } b \neq 0 \\
    1 & \text{if every column has a pivot (no free variables)} \\
    \infty & \text{if there are free variables and the system is consistent}
\end{cases}
\]

\[
\text{rank}(A) + \text{nullity}(A) = n \quad \text{(Rank--Nullity Theorem)}
\]

\subsection{Practice Questions}

\begin{enumerate}
    \item Solve the system:
    \[
        \begin{cases}
            2x + y - z = 8 \\
            -3x - y + 2z = -11 \\
            -2x + y + 2z = -3
        \end{cases}
    \]
    \textbf{Hint:} Form the augmented matrix and row reduce.

    \item If $A$ is a $5 \times 7$ matrix with $\operatorname{rank}(A) = 4$, what is the dimension of the null space of $A$?

    \item Under what conditions is $A\mathbf{x} = \mathbf{0}$ guaranteed to have only the trivial solution?
\end{enumerate}

\section{Topic 2: Vector Spaces}

\subsection{Key Concepts}

A vector space $V$ over a field $F$ is a set equipped with vector addition and scalar multiplication satisfying 10 axioms (closure, associativity, commutativity, identity, inverses, and distributivity).

\begin{itemize}
    \item \textbf{Subspace:} A subset $W \subseteq V$ that is closed under addition and scalar multiplication and contains $\mathbf{0}$.
    \item \textbf{Span:} $\operatorname{span}(\mathbf{v}_1, \ldots, \mathbf{v}_k) = \{c_1\mathbf{v}_1 + \cdots + c_k\mathbf{v}_k \mid c_i \in F\}$.
    \item \textbf{Linear independence:} Vectors $\{\mathbf{v}_1, \ldots, \mathbf{v}_k\}$ are linearly independent if $c_1\mathbf{v}_1 + \cdots + c_k\mathbf{v}_k = \mathbf{0}$ implies $c_1 = \cdots = c_k = 0$.
    \item \textbf{Basis:} A linearly independent spanning set.
    \item \textbf{Dimension:} The number of vectors in any basis.
\end{itemize}

\subsection{Key Formulas}

\[
\dim(V) = \dim(\operatorname{range}(T)) + \dim(\operatorname{null}(T))
\]

For the column space: $\operatorname{Col}(A) = \operatorname{span}$ of the pivot columns of $A$.

For the null space: $\operatorname{Null}(A) = \{\mathbf{x} \mid A\mathbf{x} = \mathbf{0}\}$.

\subsection{Practice Questions}

\begin{enumerate}
    \item Prove that the set of all $2 \times 2$ symmetric matrices forms a subspace of $M_{2 \times 2}(\mathbb{R})$.
    \item Find a basis for the subspace of $\mathbb{R}^4$ spanned by $\{(1,2,3,4), (2,1,3,4), (1,1,1,1)\}$.
    \item What is the dimension of the space of polynomials of degree at most $n$?
\end{enumerate}

\section{Topic 3: Eigenvalues and Diagonalization}

\subsection{Key Concepts}

\begin{itemize}
    \item \textbf{Eigenvalue/eigenvector:} $\lambda$ is an eigenvalue of $A$ with eigenvector $\mathbf{v}$ if $A\mathbf{v} = \lambda\mathbf{v}$, where $\mathbf{v} \neq \mathbf{0}$.
    \item \textbf{Characteristic polynomial:} $p(\lambda) = \det(A - \lambda I)$.
    \item \textbf{Algebraic multiplicity:} The multiplicity of $\lambda$ as a root of $p(\lambda)$.
    \item \textbf{Geometric multiplicity:} $\dim(\operatorname{Null}(A - \lambda I))$, the number of linearly independent eigenvectors for $\lambda$.
    \item \textbf{Diagonalizable:} $A$ is diagonalizable if and only if $A$ has $n$ linearly independent eigenvectors, or equivalently, if the sum of geometric multiplicities equals $n$.
\end{itemize}

\subsection{Key Formulas}

\[
\det(A - \lambda I) = 0 \quad \text{(characteristic equation)}
\]

\[
A = PDP^{-1}, \quad \text{where } D = \operatorname{diag}(\lambda_1, \ldots, \lambda_n), \; P = [\mathbf{v}_1 \mid \cdots \mid \mathbf{v}_n]
\]

\[
A^k = PD^kP^{-1} \quad \text{(used to compute powers of } A\text{)}
\]

\subsection{Practice Questions}

\begin{enumerate}
    \item Find the eigenvalues and eigenvectors of:
    \[
        A = \begin{pmatrix} 4 & 1 \\ 2 & 3 \end{pmatrix}
    \]
    \item Is the matrix $A = \begin{pmatrix} 1 & 1 \\ 0 & 1 \end{pmatrix}$ diagonalizable? Explain.
    \item If $A$ is diagonalizable with eigenvalues $2$ and $5$, what are the eigenvalues of $A^3 - 2A + I$?
\end{enumerate}

\section{Topic 4: Orthogonality and Inner Products}

\subsection{Key Concepts}

\begin{itemize}
    \item \textbf{Inner product:} A function $\langle \mathbf{u}, \mathbf{v} \rangle$ satisfying linearity, symmetry, and positive-definiteness.
    \item \textbf{Orthogonal vectors:} $\mathbf{u} \perp \mathbf{v}$ if $\langle \mathbf{u}, \mathbf{v} \rangle = 0$.
    \item \textbf{Norm:} $\|\mathbf{v}\| = \sqrt{\langle \mathbf{v}, \mathbf{v} \rangle}$.
    \item \textbf{Gram--Schmidt process:} Given $\{\mathbf{v}_1, \ldots, \mathbf{v}_k\}$, construct an orthogonal set $\{\mathbf{u}_1, \ldots, \mathbf{u}_k\}$:
\end{itemize}

\[
    \mathbf{u}_1 = \mathbf{v}_1, \qquad \mathbf{u}_j = \mathbf{v}_j - \sum_{i=1}^{j-1} \frac{\langle \mathbf{v}_j, \mathbf{u}_i \rangle}{\langle \mathbf{u}_i, \mathbf{u}_i \rangle} \mathbf{u}_i
\]

\begin{itemize}
    \item \textbf{Orthonormal basis:} An orthogonal basis in which each vector has unit norm.
    \item \textbf{Orthogonal projection:} $\operatorname{proj}_{\mathbf{u}}(\mathbf{v}) = \dfrac{\langle \mathbf{v}, \mathbf{u} \rangle}{\langle \mathbf{u}, \mathbf{u} \rangle} \mathbf{u}$.
\end{itemize}

\subsection{Practice Questions}

\begin{enumerate}
    \item Apply the Gram--Schmidt process to $\{(1,1,0), (1,0,1), (0,1,1)\}$ in $\mathbb{R}^3$.
    \item Find the projection of $\mathbf{b} = (1,2,3)$ onto the subspace spanned by $\mathbf{a} = (1,0,1)$.
    \item State the Cauchy--Schwarz inequality and explain its geometric meaning.
\end{enumerate}

\section{Study Schedule}

\begin{center}
\begin{tabular}{|l|l|l|}
\hline
\textbf{Day} & \textbf{Topic} & \textbf{Activities} \\
\hline
Mon, Mar 9 & Systems of equations & Review notes, 10 practice problems \\
\hline
Tue, Mar 10 & Vector spaces & Review notes, 8 practice problems \\
\hline
Wed, Mar 11 & Eigenvalues/diagonalization & Review notes, 10 practice problems \\
\hline
Thu, Mar 12 & Orthogonality & Review notes, 8 practice problems \\
\hline
Fri, Mar 13 & Mixed review & 1 full practice exam \\
\hline
Sat, Mar 14 & Light review & Review index card, weak areas only \\
\hline
Sun, Mar 15 & \textbf{EXAM DAY} & Arrive 15 minutes early \\
\hline
\end{tabular}
\end{center}

\section{Quick Reference: Essential Properties}

\begin{enumerate}
    \item $\operatorname{rank}(A) = \operatorname{rank}(A^T)$
    \item $\det(AB) = \det(A)\det(B)$
    \item $\lambda$ is an eigenvalue of $A^{-1}$ iff $1/\lambda$ is an eigenvalue of $A$ (for invertible $A$).
    \item A symmetric matrix is always diagonalizable and has real eigenvalues.
    \item The eigenvectors corresponding to distinct eigenvalues are linearly independent.
    \item An orthogonal matrix $Q$ satisfies $Q^TQ = QQ^T = I$ and preserves norms: $\|Q\mathbf{x}\| = \|\mathbf{x}\|$.
\end{enumerate}

\end{document}
